
\begin{proposition}[The CDP Euler sign and its complete geometric equality criterion]\label{cdp:st:C.2.G.1}
Let $Y$ be a connected smooth compact complex threefold with
$b_1(Y)=b_2(Y)=0$, and let $f:Y\to B$ be a proper surjective
holomorphic map with connected fibres and smooth general fibre a
complex two-torus. No condition on the singularities or
multiplicities of its fibres is imposed. Then $B\simeq\PP^1$ and
\[
 c_3(Y)=e(Y)\in\{0,2\},\qquad b_3(Y)=2-c_3(Y).       \tag{C.G.1}\label{cdp:eq:C.G.1}
\]
Every reduced fibre support has nonnegative Euler number.
An irreducible singular support with rational normalization
resolution has Euler number two; every other irreducible support
has Euler number zero.

The equality $c_3(Y)=0$ holds precisely when every fibre support
has one of the following geometric descriptions. An irreducible
smooth support has a finite unramified torus cover. An irreducible
singular support has smooth minimal elliptic ruled normalization,
or has normalization resolution of class VII. A reducible support has
normalization a disjoint union of smooth minimal Hopf or
non-K\"ahler elliptic surfaces; its reduced conductor image and
inverse image are disjoint unions of smooth elliptic curves,
and the maps between these curves are finite unramified covers.
These descriptions concern the actual normalization and
resolution maps of the given fibres; all conductor multiplicities
and full-fibre multiplicities are retained.

In the ruled case the geometric alternatives in (C.IZ.1) are
exhaustive. A contracted horizontal section is excluded by
C.2.UN.7: its fundamental group would be trivial, whereas
\ref{cdp:st:D.1.9} supplies a nonzero integral degree-one class on this
support. The class VII data necessarily satisfy the exact
equalities (C.VII.9), sharpened to (C.IZ.2). This criterion
does not assert that an arbitrary abstract collection of such
surface and curve data admits a torus smoothing.
\end{proposition}
\begin{proof}
The injection $H^1(B,\Q)\to H^1(Y,\Q)$ in the connected-fibre
Leray sequence gives $g(B)=0$. We may therefore apply the local
transgression and finite-character constructions over $\PP^1$.

First let the support $S$ be irreducible. A normal singular
support is excluded by \ref{cdp:st:A.5.2}. A smooth support has Euler number
zero. More explicitly, for a full fibre $mS$ the normalized base
change $t=s^m$ is the finite unramified map of A.3.4: in a local
equation $t=uz^m$ it is given by $w^m=u$ and $s=wz$.
Its central reduction is smooth, and its nearby fibres are the
original tori. Proper smooth deformation identifies their
integral cohomology and transports the first Chern class of
the relative tangent bundle, so the central cover has $b_1=4$
and rational first Chern class zero. A compact complex surface
with even $b_1$ is K\"ahler. In the surface classification a
K\"ahler surface with rational first Chern class zero has
Kodaira dimension zero: a nonminimal model has a $(-1)$-curve
pairing nontrivially with its canonical class; a minimal
rational or ruled surface has a line or ruling fibre with
negative canonical degree; and positive Kodaira dimension
would give a nonzero
effective pluricanonical divisor pairing positively with a
K\"ahler class. The Kodaira-dimension-zero list and $b_1=4$
therefore identify the cover as a torus. One may also compute
$e(S)=0$ directly: the
Cartier-divisor Riemann--Roch calculation gives
$\chi(\OO_S)=0$, and adjunction gives $c_1(K_S)=0$ rationally,
so Noether's formula gives $e(S)=0$.

For a nonnormal irreducible support let $Z\to T\to S$ be its
minimal normalization resolution. The effective nonzero
resolution conductor $D$ satisfies $K_Z+D$ torsion by \ref{cdp:st:A.5.3}
and the canonical evaluation near this full fibre. If any
positive tensor power of $K_Z$ had a nonzero section, taking
a further power killing that torsion would give a nonzero
section of $\OO_Z(-aD)$ for some $a>0$. On a compact connected
surface such a section would be a holomorphic function
vanishing along a nonempty divisor, which is impossible.
Thus $\kappa(Z)=-\infty$. The surface classification gives a
rational surface, an irrational ruled surface, or a class VII
surface, with the actual blowup and normalization morphisms
retained.

For rational $Z$, \ref{cdp:st:C.2.MH.2} gives $e(S)=2$. For irrational
ruled $Z$, the reduced-horizontal case is \ref{cdp:st:C.2.IRR.1}. The
coefficient-two horizontal cases are the cusp and tacnode
comparisons \ref{cdp:st:A.5.18} and \ref{cdp:st:A.5.20} and the complete ordinary-triple
comparisons \ref{cdp:st:A.5.19} and \ref{cdp:st:C.2.D4.1}--2. When the horizontal curve
is not contracted, these prove that the ruling base is
elliptic and the normalization is smooth and minimal.
A contracted section cannot occur: \ref{cdp:st:C.2.UN.7} gives
$\pi_1(S)=1$ in that case, contradicting the nonzero
restriction of the integral class of \ref{cdp:st:D.1.9}.
We make the exhaustion explicit. For a noncontracted $2C$,
degree one over the normalized image makes that image have
positive genus; no vertical rational curve can dominate it,
so conductor exponent two is the single-branch ordinary
cusp. In degree two the two horizontal branches each have
conductor exponent two. They are smooth, and their total
intersection contributions allow precisely contact two
between them, or simple intersections with one additional
smooth branch: the tacnode and ordinary triple types.
Degree three permits precisely three transverse smooth
branches. Degree at least four would already give
conductor exponent at least three on each branch, since
each meets the other three branches, and is impossible.
Together with the reduced-horizontal case and the exclusion
of a contracted section,
these exhaust the horizontal degree two of the
anticanonical divisor on a ruling fibre. \ref{cdp:st:C.2.MH.2} now gives
$e(S)=0$ for all of them. Finally \ref{cdp:st:C.2.MH.4} gives $e(S)=0$ for
class VII $Z$, retaining the exact equality data proved in
\ref{cdp:st:C.2.VII.4}.

Every reducible support has nonnegative Euler number by
\ref{cdp:st:C.2.INT.3}. That theorem gives both directions of precisely
the normalization, elliptic-conductor and unramified-cover
description stated here. Its proof works on the whole
reduced support and retains the unequal multiplicities of
the original fibre, so it applies to the remaining fibres
without a base-change substitution.

The critical values are finite and the map is a smooth
locally trivial bundle away from them, as proved in \ref{cdp:st:C2a}.
The Euler number of its torus fibre is zero, so that theorem
and complex-oriented Gauss--Bonnet give
\[
                  c_3(Y)=e(Y)=\sum_{b\in B}e((Y_b)_{\rm red})\ge0.
                                                        \tag{C.G.2}\label{cdp:eq:C.G.2}
\]
Poincar\'e duality on the smooth oriented six-manifold gives
$e(Y)=2-b_3(Y)$. Its nonsingular middle pairing is
skew-symmetric, so $b_3(Y)$ is an even nonnegative integer.
Together these facts give \eqref{cdp:eq:C.G.1}. All summands in \eqref{cdp:eq:C.G.2}
are nonnegative; their sum is zero exactly when every
summand is zero. The irreducible calculation and the
reducible equivalence give the stated complete criterion.
\end{proof}


